\documentclass[a4paper,11pt]{article}

% ---------- Language and encoding ----------
\usepackage[utf8]{inputenc}
\usepackage[T1]{fontenc}
\usepackage[english]{babel}

% ---------- Layout ----------
\usepackage[margin=2.5cm]{geometry}
\usepackage{fancyhdr}
\pagestyle{fancy}
\fancyhf{}
\lhead{Integrated Systems Technology}
\rhead{Laboratory 1}
\cfoot{\thepage}

% ---------- Math and units ----------
\usepackage{amsmath,amssymb}
\usepackage{siunitx}

% ---------- Figures, tables ----------
\usepackage{graphicx}
\usepackage{float}
\usepackage{booktabs}
\usepackage{caption}
\usepackage{subcaption}
\captionsetup{font=small, labelfont=bf}
\graphicspath{{images/}}   % put screenshots in the "images" folder

% ---------- Code ----------
\usepackage{listings}
\usepackage{xcolor}
\lstset{basicstyle=\ttfamily\footnotesize, breaklines=true, frame=single,
        keywordstyle=\color{blue}, commentstyle=\color{gray}}

% ---------- Helpers ----------
% Visible reminder of what is still to be written
\newcommand{\todo}[1]{\textcolor{red}{[TODO: #1]}}
% Placeholder box: replace with \includegraphics when the screenshot is ready
\newcommand{\figplaceholder}[1]{%
  \fbox{\parbox{0.8\textwidth}{\centering\vspace{1.5cm}#1\vspace{1.5cm}}}}

\usepackage[hidelinks]{hyperref}

% ---------- Report data: edit here ----------
\title{\textbf{Integrated Systems Technology}\\[0.5em]
       Laboratory 1: Bulk MOSFET process and electrical characterization}
\author{Name Surname \and Name Surname \\ Group N}
\date{\today}

\begin{document}

\maketitle
\tableofcontents
\newpage

% =====================================================
\section{Introduction}
Aim of the laboratory: simulate with Synopsys Sentaurus the fabrication process
of an n-type bulk MOSFET with nominal gate length $L = \SI{0.18}{\micro\meter}$
(Sentaurus Process) and then extract its electrical characteristics
(Sentaurus Device, Sentaurus Inspect).
\todo{Short description of the workflow and of the tools used: Process, Device, Visual, Inspect.}

% =====================================================
\section{Process simulation}

\subsection{Substrate definition and initial mesh}
Initial domain, mesh strategy (fine at the surface, coarser in depth) and
n-type substrate with $N_D = \SI{e15}{\per\cubic\centi\meter}$ of phosphorus.
\todo{Comment on the mesh and why it is graded; optional snapshot of Fig.~1.1 equivalent.}

\begin{figure}[H]
  \centering
  \figplaceholder{SVISUAL: initial mesh of the substrate}
  \caption{Initial 2D mesh of the substrate.}
  \label{fig:mesh_initial}
\end{figure}

\subsection{Boron implantation (p-well, punch-through stop, $V_{th}$ adjust)}
Three boron implants (\SI{200}{\kilo\electronvolt}, \SI{80}{\kilo\electronvolt},
\SI{25}{\kilo\electronvolt}) followed by an anneal at \SI{1050}{\celsius} for
\SI{10}{\second}.

\begin{figure}[H]
  \centering
  \figplaceholder{Inspect: B profiles after each step (.plx files)}
  \caption{Boron concentration versus depth after each implant and after the anneal.}
  \label{fig:boron_profiles}
\end{figure}

\todo{Zoom in, explain the differences between the profiles (peak position and
height versus energy and dose, effect of the anneal), supported by the .tdr snapshots.}

\subsection{Gate oxide growth}
Dry oxidation at \SI{850}{\celsius} for \SI{10}{\minute}.

\begin{figure}[H]
  \centering
  \figplaceholder{SVISUAL: oxide on the substrate, with ruler measurement}
  \caption{Gate oxide after thermal oxidation.}
  \label{fig:oxide}
\end{figure}

\todo{Oxide thickness measured with the ruler: $t_{ox} = \ldots$ nm.}

\subsection{Gate creation and polysilicon reoxidation}
Polysilicon deposition (\SI{0.18}{\micro\meter}), patterning with \texttt{gate\_mask}
(half transistor, mirrored at the end), anisotropic etch of polysilicon and oxide,
reoxidation at \SI{900}{\celsius} for \SI{10}{\minute}.
\todo{Snapshots of NMOS1, NMOS2, NMOS3 and a brief comment on each etch step.}

\subsection{LDD implantation}
Arsenic implant (\SI{4e14}{\per\square\centi\meter}, \SI{10}{\kilo\electronvolt})
and short anneal. Refinement box defined before the implant.
\todo{Snapshot of the activated As profile (AsActive) with a suitable range.}

\subsection{HALO implantation}
Four boron implants with tilt \SI{30}{\degree} and rotations
\SI{0}{\degree}, \SI{90}{\degree}, \SI{180}{\degree}, \SI{270}{\degree},
followed by an RTA of \SI{5}{\second} at \SI{1050}{\celsius}.

\begin{figure}[H]
  \centering
  \begin{subfigure}{0.48\textwidth}
    \centering\figplaceholder{Halo 1}
    \caption{Halo 1 (rotation \SI{0}{\degree}).}
  \end{subfigure}\hfill
  \begin{subfigure}{0.48\textwidth}
    \centering\figplaceholder{Halo 2}
    \caption{Halo 2 (rotation \SI{90}{\degree}).}
  \end{subfigure}

  \vspace{1em}
  \begin{subfigure}{0.48\textwidth}
    \centering\figplaceholder{Halo 3}
    \caption{Halo 3 (rotation \SI{180}{\degree}).}
  \end{subfigure}\hfill
  \begin{subfigure}{0.48\textwidth}
    \centering\figplaceholder{Halo 4}
    \caption{Halo 4 (rotation \SI{270}{\degree}).}
  \end{subfigure}
  \caption{Structure after each HALO implantation.}
  \label{fig:halo}
\end{figure}

\begin{figure}[H]
  \centering
  \figplaceholder{SVISUAL: structure after Halo RTA (NMOS10)}
  \caption{Boron profile after the HALO RTA.}
  \label{fig:halo_rta}
\end{figure}

\todo{Comment the effect of each of the four implants and of the RTA (symmetry,
reduction of charge sharing / punch-through).}

\subsection{Spacers formation}
Isotropic nitride deposition (\SI{60}{\nano\meter}), anisotropic etch with a small
isotropic overetch, removal of the thin oxide.
\todo{Snapshots of NMOS11, NMOS12, NMOS13 and a brief comment.}

\subsection{Source/drain implantation}
Refinement box in the S/D region (coordinates of Fig.~1.11 of the handout, uniform
spacing of \SI{10}{\nano\meter} along $x$ and \SI{50}{\nano\meter} along $y$).

\begin{lstlisting}
% paste here the refinebox command lines actually used
refinebox Silicon min= {...} max= {...} xrefine= {...} yrefine= {...} add
grid remesh
\end{lstlisting}

\todo{Comment the command lines (meaning of min, max, xrefine, yrefine, grid remesh).}

\begin{figure}[H]
  \centering
  \figplaceholder{SVISUAL zoom with mesh enabled}
  \caption{Local mesh refinement in the S/D region.}
  \label{fig:refinebox_sd}
\end{figure}

Then Arsenic implant (\SI{5e15}{\per\square\centi\meter}, \SI{40}{\kilo\electronvolt},
tilt \SI{7}{\degree}) and final RTA.
\todo{Snapshot of NMOS14 and NMOS15.}

\subsection{Contacts and final structure}
Aluminium deposition and patterning (S/D pads and gate pad), mirroring of the
half structure (\texttt{transform reflect left}), definition of the contacts
(substrate, gate, source, drain).
\todo{Snapshots of NMOS16--NMOS18, of the final structure and of the presimulation structure.}

% =====================================================
\section{Electrical simulation}

\subsection{Simulation setup}
Brief description of the \texttt{sdevice\_des.cmd} file: sections File, Electrode,
Physics (doping-dependent mobility, high-field saturation, SRH), Math and Solve
(ramp of $V_{DS}$ up to \SI{2}{\volt}, then ramp of $V_{GS}$ up to \SI{2}{\volt}).
\todo{Add the two additional quasistationary blocks for $V_{DS} = \SI{0.2}{\volt}$ and
the \texttt{Load} / \texttt{NewCurrentPrefix} lines used.}

\subsection{Transcharacteristic}
\begin{figure}[H]
  \centering
  \figplaceholder{Inspect: $I_{DS}(V_{GS})$, log scale, $V_{DS} = \SI{2}{\volt}$}
  \caption{$I_{DS}(V_{GS})$ characteristic in semi-logarithmic scale at $V_{DS} = \SI{2}{\volt}$.}
  \label{fig:idvg_high}
\end{figure}

\subsection{$I_{OFF}$ and $I_{ON}$}
\begin{table}[H]
  \centering
  \caption{Off and on currents ($V_{DS} = \SI{2}{\volt}$, per \si{\micro\meter} of width).}
  \label{tab:ion_ioff}
  \begin{tabular}{lcc}
    \toprule
    Quantity & Definition & Value \\
    \midrule
    $I_{OFF}$ & $I_{DS}$ at $V_{GS} = \SI{0}{\volt}$ & \todo{value} \si{\ampere\per\micro\meter} \\
    $I_{ON}$  & $I_{DS}$ at $V_{GS} = \SI{2}{\volt}$ & \todo{value} \si{\ampere\per\micro\meter} \\
    \bottomrule
  \end{tabular}
\end{table}

\subsection{Threshold voltage and subthreshold slope}
$V_{th}$ from the second-derivative method and SS from the maximum slope of
$\log_{10} I_{DS}$:
\begin{equation}
  V_{th} = V_{GS}\Big|_{\max \frac{d^2 I_{DS}}{dV_{GS}^2}}, \qquad
  SS = \left[\max \frac{d\,\log_{10} I_{DS}}{dV_{GS}}\right]^{-1}
\end{equation}

\begin{table}[H]
  \centering
  \caption{Figures of merit extracted with Sentaurus Inspect.}
  \label{tab:fom}
  \begin{tabular}{lcc}
    \toprule
    & $V_{DS} = \SI{0.2}{\volt}$ & $V_{DS} = \SI{2}{\volt}$ \\
    \midrule
    $V_{th}$ (\si{\volt})        & \todo{} & \todo{} \\
    $SS$ (\si{\milli\volt/dec})  & \todo{} & \todo{} \\
    \bottomrule
  \end{tabular}
\end{table}

\subsection{DIBL}
\begin{equation}
  \mathrm{DIBL} = \frac{V_{th}(V_{DS,low}) - V_{th}(V_{DS,high})}{V_{DS,high} - V_{DS,low}}
  \quad [\si{\milli\volt/\volt}]
\end{equation}
\todo{Plot the two transcharacteristics on the same graph, compute DIBL and comment the value.}

\begin{figure}[H]
  \centering
  \figplaceholder{Inspect: $I_{DS}(V_{GS})$ at $V_{DS} = \SI{0.2}{\volt}$ and \SI{2}{\volt}}
  \caption{Transcharacteristics at low and high drain bias.}
  \label{fig:idvg_both}
\end{figure}

% =====================================================
\section{Conclusions}
\todo{Summarize the main results ($V_{th}$, SS, DIBL, $I_{ON}/I_{OFF}$), comment on the
role of LDD, HALO and punch-through implants, and on short-channel effects of the
\SI{0.18}{\micro\meter} device.}

\end{document}